\chapter{Regularization of the Master Problem}
\label{chap:regularization}

\vspace*{\fill}

This chapter replaces the linear restricted master problem by a regularized one. The
regularization is formulated in the dual space, where it adds a quadratic penalty on the
distance between the dual iterate and a \textbf{stability center}. Translated back to the
primal space, it yields a convex quadratic program with an auxiliary variable, the
\textbf{explicit regularized master problem}.

\medskip

Because all linking rows of the unified formulation are equalities, the auxiliary variable
can be eliminated for equality and inequality rows alike. The result, the
\textbf{eliminated regularized master problem}, is a convex quadratic program whose only
constraints are the convexity constraints and the nonnegativity of the variables. This is the
problem that the active-set method of the following chapters solves.

\medskip

Two identities derived in this chapter are used throughout the rest of the work. The first
expresses the duals of the linking rows as an explicit function of the primal iterate. The
second shows that the gradient of the quadratic objective consists of the reduced costs of
the columns and the duals of the inequality rows. The chapter ends with the regularized
column generation algorithm, its exact stopping criterion and the accuracy that this
criterion guarantees.

\vspace*{\fill}

\newpage


\section{Proximal Stabilization of the Dual}
\label{sec:proximal}

Consider the restricted master problem in the unified form~\eqref{eq:unified}, over the
columns $\mathcal{J}'$ and all $m_I$ surpluses. Its dual is
\begin{equation} \label{eq:restricted-dual}
\max_{\pi,\,\mu}\; b^\top\pi + \one^\top\mu
\qquad\text{s.t.}\qquad
c_j - a_j^\top\pi - \mu_{\psi(j)}\ge0 \;\;(j\in\mathcal{J}'),\qquad \pi_I\ge0 .
\end{equation}
Its optimal dual solutions are extreme points of the feasible region, which causes the
instability described in \cref{rem:instability}. Proximal stabilization adds a penalty on the
distance to a given point $\pihat\in\R^m$.

\begin{definition}[Stabilized dual problem] \label{def:stabilized-dual}
Let $\pihat\in\R^m$ be the \textbf{stability center} and $\delta>0$ the
\textbf{stabilization parameter}. The \textbf{stabilized dual problem} is
\begin{equation} \label{eq:stabilized-dual}
\max_{\pi,\,\mu}\; b^\top\pi + \one^\top\mu - \frac{\delta}{2}\,\lVert\pi-\pihat\rVert^2
\qquad\text{s.t.}\qquad
c_j - a_j^\top\pi - \mu_{\psi(j)}\ge0 \;\;(j\in\mathcal{J}'),\qquad \pi_I\ge0 .
\end{equation}
\end{definition}

The objective of~\eqref{eq:stabilized-dual} is strictly concave in $\pi$, so the optimal
$\pi$ is unique. A large $\delta$ keeps $\pi$ close to $\pihat$ (strong stabilization), and
$\delta\to0$ recovers~\eqref{eq:restricted-dual}. Replacing the center $\pihat$ by the
solution $\pi$ and solving again is one step of the \textbf{proximal point method}
\cite{rockafellar1976monotone} applied to the restricted dual problem. If the center is only
moved when it improves the Lagrangian bound, we obtain a \textbf{proximal bundle method}
\cite{kiwiel1990proximity, briant2008comparison}. Both rules are discussed in
\cref{sec:center-update}.


\section{The Explicit Regularized Master Problem}
\label{sec:explicit}

The primal counterpart of~\eqref{eq:stabilized-dual} has an auxiliary vector
$w\in\R^m$, one entry per linking row.

\begin{definition}[Explicit regularized master problem] \label{def:explicit}
The \textbf{explicit regularized master problem} is
\begin{subequations} \label{eq:explicit}
\begin{align}
\min_{y,\,w}\quad & \cbar^\top y + \frac{\delta}{2}\lVert w + \pihat\rVert^2
                     - \frac{\delta}{2}\lVert\pihat\rVert^2 \\
\text{s.t.}\quad & \Abar y + \delta w = b, && (\pi) \label{eq:explicit-link}\\
& \Cbar y = \one, && (\mu)\\
& y\ge0 .
\end{align}
\end{subequations}
\end{definition}

\begin{theorem} \label{thm:explicit-duality}
The dual of~\eqref{eq:explicit} is the stabilized dual problem~\eqref{eq:stabilized-dual}.
Strong duality holds, and at a pair of optimal solutions
\begin{equation} \label{eq:pi-w}
\pi = \pihat + w .
\end{equation}
\end{theorem}

\begin{proof}
The Lagrangian of~\eqref{eq:explicit}, with multipliers $\pi$, $\mu$ and $\nu\ge0$ for
$y\ge0$, is
\[
\mathcal{L} = \cbar^\top y + \tfrac{\delta}{2}\lVert w+\pihat\rVert^2
- \tfrac{\delta}{2}\lVert\pihat\rVert^2
+ \pi^\top(b - \Abar y - \delta w) + \mu^\top(\one - \Cbar y) - \nu^\top y .
\]
It is a strictly convex quadratic function of $w$, minimized where
$\delta(w + \pihat) - \delta\pi = 0$, i.e.\ at $w = \pi - \pihat$. Inserting this,
\[
\tfrac{\delta}{2}\lVert\pi\rVert^2 - \tfrac{\delta}{2}\lVert\pihat\rVert^2
- \delta\pi^\top(\pi-\pihat) = -\tfrac{\delta}{2}\lVert\pi-\pihat\rVert^2 .
\]
The Lagrangian is linear in $y$, and its infimum over $y$ is finite only if
$\cbar - \Abar^\top\pi - \Cbar^\top\mu - \nu = 0$ for some $\nu\ge0$, i.e.\ only if
$\cbar - \Abar^\top\pi - \Cbar^\top\mu\ge0$. The dual function is therefore
$b^\top\pi + \one^\top\mu - \frac{\delta}{2}\lVert\pi-\pihat\rVert^2$, subject to this
condition. Written per column, the condition for the columns of $\lambda$
is~\eqref{eq:dual-columns}, and the condition for the surplus columns is $\pi_I\ge0$
(\cref{lem:unified}). This is~\eqref{eq:stabilized-dual}. Problem~\eqref{eq:explicit} is a
feasible convex quadratic program with linear constraints (it is feasible because $w$ is
free), so strong duality holds, and the minimizing $w$ satisfies~\eqref{eq:pi-w}. 
\end{proof}

\begin{remark} \label{rem:sign-in-dual}
\Cref{thm:explicit-duality} shows that the sign condition $\pi_I\ge0$ is part of the
stabilized dual problem automatically, through the surplus columns. The stability center
$\pihat$ itself may violate it. The proximal step then moves $\pi$ back into the region
$\pi_I\ge0$.
\end{remark}

\begin{remark} \label{rem:why-surplus}
Without surplus variables, an inequality row $(A_I\lambda)_i\ge(b_I)_i$ would enter the
explicit problem as $(A_I\lambda)_i + \delta w_i\ge(b_I)_i$, and $w_i$ would no longer be
determined by $\lambda$. The surpluses make every row of~\eqref{eq:explicit-link} an
equality, so that $w$ can be eliminated for all rows in the next section.
\end{remark}


\section{The Eliminated Regularized Master Problem}
\label{sec:eliminated}

From~\eqref{eq:explicit-link} we get $w = (b - \Abar y)/\delta$ for every feasible point.
Substituting this into the objective of~\eqref{eq:explicit} removes $w$ completely.

\begin{theorem}[Eliminated regularized master problem] \label{thm:eliminated}
Problem~\eqref{eq:explicit} is equivalent to the convex quadratic program
\begin{subequations} \label{eq:eliminated}
\begin{align}
\min_{y}\quad & f(y) = \frac{1}{2\delta}\, y^\top\Abar^\top\Abar\, y + \ctil^\top y \\
\text{s.t.}\quad & \Cbar y = \one, \label{eq:eliminated-conv}\\
& y\ge0, \label{eq:eliminated-bounds}
\end{align}
\end{subequations}
with the linear cost
\begin{equation} \label{eq:ctilde}
\ctil = \cbar - \Abar^\top\Bigl(\pihat + \frac{b}{\delta}\Bigr).
\end{equation}
More precisely, $(y, w)$ is optimal for~\eqref{eq:explicit} if and only if $y$ is optimal
for~\eqref{eq:eliminated} and $w = (b - \Abar y)/\delta$. The objective values are related by
\begin{equation} \label{eq:objective-relation}
\Phi(y) := \cbar^\top y + \frac{\delta}{2}\lVert w+\pihat\rVert^2 - \frac{\delta}{2}\lVert\pihat\rVert^2
= f(y) + \frac{1}{2\delta}\lVert b\rVert^2 + \pihat^\top b .
\end{equation}
\end{theorem}

\begin{proof}
Let $r = b - \Abar y$, so $w = r/\delta$. Then
\[
\tfrac{\delta}{2}\lVert r/\delta + \pihat\rVert^2 - \tfrac{\delta}{2}\lVert\pihat\rVert^2
= \tfrac{1}{2\delta}\lVert r\rVert^2 + \pihat^\top r
= \tfrac{1}{2\delta}\lVert\Abar y\rVert^2 - \tfrac{1}{\delta}b^\top\Abar y
  + \tfrac{1}{2\delta}\lVert b\rVert^2 + \pihat^\top b - \pihat^\top\Abar y .
\]
Adding $\cbar^\top y$ and collecting the terms that are linear in $y$ gives $f(y)$ plus the
constant in~\eqref{eq:objective-relation}. The constraint~\eqref{eq:explicit-link} only
defines $w$, so the remaining constraints are~\eqref{eq:eliminated-conv}
and~\eqref{eq:eliminated-bounds}. 
\end{proof}

The Hessian of $f$ is $\frac{1}{\delta}\Abar^\top\Abar$, with the block structure
\begin{equation} \label{eq:hessian-blocks}
\frac{1}{\delta}\Abar^\top\Abar
= \frac{1}{\delta}\begin{bmatrix} A_E^\top A_E + A_I^\top A_I & -A_I^\top\\ -A_I & I\end{bmatrix}.
\end{equation}
It is positive semidefinite but in general singular: it has rank at most $m$, while the
number of variables $n + m_I$ grows with every generated column. Problem~\eqref{eq:eliminated}
nevertheless has an optimal solution, because its feasible region is bounded in the $\lambda$
part and the objective grows quadratically in $s$.

\subsection{The dual vector as a function of the primal iterate}

Relation~\eqref{eq:pi-w} with $w = (b-\Abar y)/\delta$ suggests the following definition,
which applies to \emph{every} $y$, not only to the optimal one.

\begin{definition}[Dual estimate] \label{def:dual-estimate}
For $y = [\lambda; s]$ the \textbf{dual estimate} is
\begin{equation} \label{eq:pi-of-y}
\pi(y) = \pihat + \frac{1}{\delta}\bigl(b - \Abar y\bigr),
\qquad\text{i.e.}\qquad
\begin{aligned}
\pi_E(y) &= \pihat_E + \tfrac{1}{\delta}\bigl(b_E - A_E\lambda\bigr),\\
\pi_I(y) &= \pihat_I + \tfrac{1}{\delta}\bigl(b_I - A_I\lambda + s\bigr).
\end{aligned}
\end{equation}
\end{definition}

A surplus raises the dual of its row: an over-covered inequality row with a positive
surplus has a larger $\pi_i$ than it would have without the surplus. The next lemma is the
central identity of this work.

\begin{lemma}[Gradient--dual relation] \label{lem:gradient-dual}
For every $y$,
\begin{equation} \label{eq:gradient-dual}
\nabla f(y) = \cbar - \Abar^\top\pi(y),
\qquad\text{i.e.}\qquad
g_j := \frac{\partial f}{\partial\lambda_j} = c_j - a_j^\top\pi(y), \qquad
\frac{\partial f}{\partial s_i} = \pi_i(y)\quad (i\in\{1,\dots,m_I\}).
\end{equation}
\end{lemma}

\begin{proof}
$\nabla f(y) = \frac{1}{\delta}\Abar^\top\Abar y + \cbar - \Abar^\top\pihat - \frac{1}{\delta}\Abar^\top b
= \cbar - \Abar^\top\bigl(\pihat + \frac{1}{\delta}(b - \Abar y)\bigr)$. The rows of
$\Abar^\top$ belonging to $\lambda_j$ are $a_j^\top$, and the rows belonging to $s_i$ are
$-e_i^\top$ (in the inequality part). 
\end{proof}

\Cref{lem:gradient-dual} has two consequences that are used throughout this work.
\begin{itemize}
    \item The gradient component of a column, $g_j = c_j - a_j^\top\pi$, is the
          \textbf{reduced cost of column $j$ without the convexity dual}. The gradient
          component of a surplus is the dual of its row.
    \item The gradient, an $(n+m_I)$-vector whose length grows with every column, is
          determined by the $m$-vector $\pi(y)$. An implementation can store $\pi$ and
          evaluate $g_j$ only where it is needed.
\end{itemize}

\begin{prop}[Objective value] \label{prop:objective-value}
For every $y$, with $\pi = \pi(y)$,
\[
\Phi(y) = c^\top\lambda + \frac{\delta}{2}\lVert\pi\rVert^2 - \frac{\delta}{2}\lVert\pihat\rVert^2 .
\]
\end{prop}

\begin{proof}
In~\eqref{eq:objective-relation}, $\Phi(y) = \cbar^\top y + \frac{\delta}{2}\lVert w+\pihat\rVert^2
- \frac{\delta}{2}\lVert\pihat\rVert^2$ with $w+\pihat = \pi(y)$, and $\cbar^\top y = c^\top\lambda$
because the surpluses have zero cost. 
\end{proof}

\begin{prop}[Uniqueness of the dual] \label{prop:pi-unique}
The optimal value of $\Abar y$, and hence the optimal dual vector $\pi$, is the same for all
optimal solutions of~\eqref{eq:eliminated}, even if the optimal $y$ is not unique.
\end{prop}

\begin{proof}
Let $y_1, y_2$ be optimal with $\Abar y_1\ne\Abar y_2$. Their midpoint $\bar y$ is feasible.
The linear part of $f$ is affine, while $\lVert\Abar y\rVert^2$ is strictly convex as a
function of $\Abar y$, so $f(\bar y) < \frac12\bigl(f(y_1)+f(y_2)\bigr)$, a contradiction.

\end{proof}

\begin{prop}[Residual and dual distance] \label{prop:residual}
For every $y$,
\[
b - \Abar y = \delta\bigl(\pi(y) - \pihat\bigr).
\]
In particular, $y$ is feasible for the unified master problem~\eqref{eq:unified} if and only
if $\pi(y) = \pihat$.
\end{prop}

\begin{proof}
This is~\eqref{eq:pi-of-y} rearranged. 
\end{proof}


\section{Regularized Column Generation}
\label{sec:center-update}

Column generation with the regularized master problem follows the classical scheme of
\cref{alg:classical-cg}, with two differences: the duals passed to the pricing problem are
the stabilized duals $\pi(y)$, and the stability center has to be updated.

\subsection{Center updates}

Two rules for the center update are common.
\begin{itemize}
    \item \textbf{Proximal point rule:} after every master problem, $\pihat\gets\pi$. Every
          iteration is a step of the proximal point method on the current restricted dual
          problem.
    \item \textbf{Bundle rule:} $\pihat\gets\pi$ only if the Lagrangian bound improves,
          $\theta(\pi) > \theta(\pihat)$ (a \emph{serious step}). Otherwise the center stays
          (a \emph{null step}). This requires the Lagrangian bound at every iteration, i.e.\
          all blocks must be priced every iteration.
\end{itemize}

\subsection{Exact termination}

In classical column generation, the pricing test alone decides optimality
(\cref{prop:cg-optimality}). With regularization this is no longer true: when pricing finds
no improving column, the regularized master problem may still be infeasible for the
unregularized one, since $\Abar y\ne b$ as long as $\pi\ne\pihat$ (\cref{prop:residual}).

\begin{theorem}[Exact termination] \label{thm:exact-termination}
Let $y = [\lambda; s]$ be optimal for the eliminated regularized master
problem~\eqref{eq:eliminated} over the columns $\mathcal{J}'$, with dual estimate
$\pi = \pi(y)$ and convexity duals $\mu$. If
\begin{enumerate}
    \item $\pi = \pihat$, and
    \item the pricing problem~\eqref{eq:pricing} has a nonnegative optimal value for every
          block,
\end{enumerate}
then $\lambda$, extended by zeros, is optimal for the master problem~\eqref{eq:master}, and
$(\pi,\mu)$ is optimal for its dual~\eqref{eq:master-dual}.
\end{theorem}

\begin{proof}
By \cref{prop:residual}, $\pi = \pihat$ gives $\Abar y = b$, i.e.\ $A_E\lambda = b_E$ and
$A_I\lambda = b_I + s\ge b_I$, so $\lambda$ is primal feasible. The optimality conditions
of~\eqref{eq:eliminated}, which are derived in \cref{thm:kkt-master} below, state that every
column of $\mathcal{J}'$ has a nonnegative reduced cost (with equality where $\lambda_j>0$),
that $\pi_I\ge0$, and that $\pi_i = 0$ whenever $s_i>0$. Condition~(ii) extends the
nonnegativity of the reduced costs to all columns. So $(\pi,\mu)$ is dual feasible, and
complementary slackness holds for the columns, the zero weights outside $\mathcal{J}'$ and the
inequality rows. 
\end{proof}

If pricing finds no improving column while $\pi\ne\pihat$, the restricted master problem
contains the right columns for the moment, but the proximal term still keeps $\pi$ away from
the restricted dual optimum. The algorithm then moves the center to $\pi$ and solves the
master problem again \emph{without adding a column}. We call this a
\textbf{center-only iteration}. It is one proximal point step on the current restricted
model. Along the way, the duals move, and pricing may find improving columns again. A run of
regularized column generation therefore consists of two kinds of iterations, which
\cref{sec:results-phases} studies in detail: iterations that add columns, and center-only
iterations that drive $\pi$ towards $\pihat$.

\begin{prop}[Accuracy of the stopping test] \label{prop:stop-accuracy}
In practice condition~(i) of \cref{thm:exact-termination} is tested as
\begin{equation} \label{eq:center-test}
\lVert\pi - \pihat\rVert_\infty \;\le\; \tau\,\max\bigl(1,\lVert\pi\rVert_\infty\bigr)
\end{equation}
with a tolerance $\tau>0$. When this test holds, the violation of the linking rows by
$\lambda$ is bounded by
\[
\lVert b - \Abar y\rVert_\infty \;\le\; \delta\,\tau\,\max\bigl(1,\lVert\pi\rVert_\infty\bigr).
\]
\end{prop}

\begin{proof}
This follows from \cref{prop:residual}. 
\end{proof}

The accuracy of the final solution is thus relative to the size of the duals, and it scales
with $\delta$. The objective value $\Phi$ and the Lagrangian bound are much less affected by
this residual than the cost $c^\top\lambda$ of a slightly infeasible $\lambda$, as the
results in \cref{sec:results-correctness} show.

\subsection{The Lagrangian bound with inequality rows}

The Lagrangian bound of \cref{prop:lagrangian-bound} requires $\pi_I\ge0$. The duals of the
regularized master problem satisfy this condition only up to the tolerance of the master
solver, and a tiny negative entry would make the bound useless (\cref{rem:bound-sign}). The
bound is therefore evaluated at the projection
\[
\pi^+ = \bigl[\,\pi_E;\;\max(\pi_I, 0)\,\bigr],
\]
which always satisfies the sign condition, so that $\theta(\pi^+)\le z^\star$ holds by
\cref{prop:lagrangian-bound}. The same projected vector is passed to the pricing problem.
For problems with equality rows only, $\pi^+ = \pi$.

\subsection{The algorithm}

\begin{algobox}{Regularized Column Generation}{alg:rcg}
\State \textbf{Input:} initial columns $\mathcal{J}'$ (one per block), center $\pihat$,
       parameter $\delta>0$, tolerances $\tau$, $\varepsilon_{\mathrm{rc}}$
\Repeat
    \State Solve the eliminated regularized master problem~\eqref{eq:eliminated}; obtain $y$,
           $\pi = \pi(y)$ and $\mu$
    \State $\pi^+\gets[\pi_E;\max(\pi_I,0)]$
    \State Price the blocks in random order with $(\pi^+,\mu)$; add the first column with
           $\bar c_j < -\varepsilon_{\mathrm{rc}}$ to $\mathcal{J}'$
    \If{no column was added (all blocks priced)}
        \State $\theta\gets\theta(\pi^+)$ \Comment{valid lower bound on $z^\star$}
        \If{$\lVert\pi-\pihat\rVert_\infty\le\tau\max(1,\lVert\pi\rVert_\infty)$}
            \State \textbf{stop}: $\lambda$ is optimal (\cref{thm:exact-termination})
        \EndIf
        \State $\pihat\gets\pi$ \Comment{center-only iteration}
    \Else
        \State $\pihat\gets\pi$ \Comment{proximal point rule (bundle rule: only if $\theta$ improves)}
    \EndIf
\Until{stopped}
\end{algobox}

In the experiments of \cref{chap:results} we use the proximal point rule,
$\delta = 1$, $\tau = 10^{-6}$ and $\varepsilon_{\mathrm{rc}} = 10^{-6}$. The remaining
chapters develop the method used to solve the master problem in the first line of the loop.
